% ============================================================
% chapter_02_literature_review.tex - Chapter 2, Literature Review
%
% Purpose: establishes the scientific context for the LSEQM+DL framework.
% Sections: 2.1 Satellite-derived Daily Precipitation, 2.2 Bias and
%           Verification, 2.3 Bias Correction Methods, 2.4 Theoretical
%           Bounds on Marginal Correction, 2.5 Reproducible Computational
%           Notebook Workflows, 2.6 Maritime Continent Regional Context.
% Included by main_thesis.tex.
% License: CC BY 4.0 (see ../LICENSE).
% ============================================================

\chapter{Literature Review}\label{chap:litreview}

This chapter establishes the scientific context for the
LSEQM+DL (Linear Scaling (LS), Empirical Quantile Mapping (EQM), Generalized Pareto Distribution (GPD), and Convolutional Neural Network (CNN) refinement, a Deep Learning (DL) architecture) framework introduced in Chapter~\ref{chap:methods}.
Section~\ref{sec:lit-imerg} reviews satellite-derived daily
precipitation products, with emphasis on the Integrated
Multi-satellitE Retrievals for Global Precipitation Measurement
(IMERG) family that the framework operates on.
Section~\ref{sec:lit-verification} reviews how bias is decomposed
into measurable dimensions and the verification standards used
to assess each dimension. Section~\ref{sec:lit-methods} reviews
the families of bias correction methods on which the four-stage
LSEQM+DL pipeline draws. Three further sections complete the
chapter: Section~\ref{sec:lit-bounds} discusses the theoretical
limits of marginal-distribution correction methods on
temporal-skill metrics; Section~\ref{sec:lit-reproducibility}
reviews the state of reproducibility in climate and hydrology
software; and Section~\ref{sec:lit-maritime-continent} reviews
the regional precipitation context of the Indonesian Maritime
Continent. Together these six sections frame the methodological
choices in Chapter~\ref{chap:methods}, the empirical findings in
Chapter~\ref{chap:results}, and the theoretical discussion in
Chapter~\ref{chap:discussion}.


% ============================================================
\section{Satellite-derived Daily Precipitation}\label{sec:lit-imerg}
% Target: 3-4 pp
% ============================================================

Daily precipitation is the basic temporal unit for many
operational hydrological and agricultural applications:
flash-flood early warning, reservoir release scheduling,
drought monitoring, crop water budgeting, and disaster-risk
insurance all consume one daily value per location per day.
For most of the global land surface, daily precipitation
fields at spatially continuous resolution are not available
from gauge networks alone. The Indonesian archipelago, the
study domain of the present work, illustrates the gap directly:
the contributing gauge network is concentrated over Java and
southern Sumatra and is sparse over Papua, Maluku, and eastern
Nusa Tenggara~\citep{ref:chen2008}. The
satellite family of precipitation products fills the gap by
providing daily fields at uniform spatial resolution across the
domain, though at the cost of systematic biases that the
present work aims to correct.

\subsection{From TRMM to GPM}

Modern satellite precipitation estimation evolved from the
Tropical Rainfall Measuring Mission (TRMM) era. TRMM operated
between 1997 and 2015 and produced the TRMM Multi-satellite
Precipitation Analysis (TMPA), which blended passive-microwave
retrievals from the TRMM platform and a constellation of
partner satellites with geostationary infrared estimates and
monthly gauge analyses~\citep{ref:huffman2007}. TMPA was the
backbone of operational satellite precipitation for nearly two
decades and seeded the verification literature that
characterised the systematic biases of microwave-plus-infrared
blends against gauges.

The Global Precipitation Measurement (GPM) mission launched in
2014 as TRMM's successor, with the Core Observatory carrying a
dual-frequency precipitation radar and a multi-channel passive
microwave imager. IMERG~\citep{ref:hou2014, ref:huffman2020} extends the
TMPA methodology to the GPM constellation and provides a
quasi-global precipitation product at $0.1^{\circ}$ spatial and
30-minute temporal resolution from 2000 to the present, with
TRMM-era retrievals reprocessed under the IMERG framework so
the product spans a continuous 25-year record.

\subsection{IMERG V07 architecture}

The IMERG algorithm~\citep{ref:huffman2020} merges three input streams.
Passive-microwave retrievals from the GPM constellation
(including the GPM Core Observatory's GMI, partner-satellite
SSMI/S, ATMS, AMSR2, MHS, and SAPHIR instruments) provide the
direct precipitation observations. The microwave samples are
sparse in time and space because each instrument visits a given
location only a few times per day; the gaps are filled by
geostationary infrared (IR) estimates that have continuous
temporal coverage but lower direct sensitivity to precipitation.
The IR estimates are merged with microwave retrievals using a
morphing scheme that propagates the microwave information
between successive overpasses, producing a continuous
half-hourly field at $0.1^{\circ}$ resolution. A final pass
anchors the monthly accumulations to the Global Precipitation
Climatology Centre (GPCC) gauge analysis, correcting long-term
mean biases while preserving the high-frequency variability of
the merged microwave-plus-IR product.

\subsection{Three latency variants}

IMERG releases three latency variants of the same daily product.
The \textit{Early Run} is released approximately 4~hours after
real time and is suitable for nowcasting applications that
require minimum delay at the cost of reduced accuracy. The
\textit{Late Run} (IMERG-L) is released approximately 14~hours
after real time and is the variant operated on in this study;
the 14-hour latency fits within the same-day decision window of
operational services such as flood early warning. The
\textit{Final Run} is released approximately 3.5 months after
real time and benefits from the monthly gauge anchoring step
that the Early and Late runs cannot include. The choice between
variants is a latency-versus-accuracy trade-off: applications
that can wait for the Final Run gain the gauge-anchored
adjustment; applications that need timely fields must accept
the larger biases of the Early and Late runs.

\subsection{Gridded daily precipitation products}

IMERG is one of several global gridded daily precipitation
products available to operational users; each makes a different
trade-off between latency, gauge anchoring, and spatial
resolution. Table~\ref{tab:gridded-products} summarises the
products most often considered for daily-scale work over the
tropics. CPC Unified Gauge-Based Analysis of Daily Precipitation (CPC-UNI) \citep{ref:chen2008} interpolates gauge
observations directly and is therefore the highest-quality
reference where the contributing network is dense, at the cost
of a coarse 0.5$^{\circ}$ grid and limited skill where the
gauge network is sparse. ERA5-Land \citep{ref:munoz2021era5land}
is a reanalysis product that uses many data streams but does
not ingest gauges directly; the gauge constraint is therefore
indirect. The IMERG Early, Late, and Final runs share the same
0.1$^{\circ}$ grid but differ in latency (about 4 hours, 14
hours, and 3.5 months respectively) and in whether the monthly
gauge anchoring step is applied (Final only). The Climate Hazards
Group InfraRed Precipitation with Stations (CHIRPS)
\citep{ref:funk2015} provides a higher-resolution gauge-anchored
product at 0.05$^{\circ}$ with two-day latency. The Precipitation
Estimation from Remotely Sensed Information using Artificial Neural
Networks Climate Data Record (PERSIANN-CDR)
\citep{ref:ashouri2015} provides a long historical record from
1983 onward, bias-adjusted to the monthly GPCP gauge analysis at
coarse resolution. The Multi-Source Weighted-Ensemble
Precipitation (MSWEP) \citep{ref:beck2019}
blends gauges, satellites, and reanalyses into a single product
at 0.1$^{\circ}$. For the present study, IMERG Late Run was
selected because its 14-hour latency fits the operational
decision window while its 0.1$^{\circ}$ resolution and 25-year
record are both adequate for the verification framework reported
in Chapter~\ref{chap:results}.

\begin{table}[H]
\centering
\caption[Daily gridded precipitation products commonly used for tropical hydrological and operational applications]{Daily gridded precipitation products commonly used for
tropical hydrological and operational applications. Latency is
the typical time lag between observation and product release;
gauge integration describes whether and how gauge data is
incorporated in the product.\label{tab:gridded-products}}
\begingroup\footnotesize
\begin{tabular}{@{}lllll@{}}
\toprule
Product & Latency & Range & Gauge integration & Resolution \\
\midrule
CPC-UNI                 & 1 day      & 1979 -- present & Full gauge analysis             & 0.5$^{\circ}$ \\
ERA5-Land               & 2--3 months & 1981 -- present & Reanalysis, no direct gauge     & $\sim$0.1$^{\circ}$ \\
IMERG Final Run         & 3 months   & 2000 -- present & Ensemble with monthly gauges    & 0.1$^{\circ}$ \\
IMERG Late Run          & 12--18 h   & 2000 -- present & No gauge integration            & 0.1$^{\circ}$ \\
IMERG Early Run         & 4--6 h     & 2000 -- present & No gauge integration            & 0.1$^{\circ}$ \\
CHIRPS                  & 2 days     & 1981 -- present & Ensemble with gauges            & 0.05$^{\circ}$ \\
PERSIANN-CDR            & 3--5 days  & 1983 -- present & Monthly GPCP adjustment         & 0.25$^{\circ}$ \\
MSWEP                   & 1--3 days  & 1979 -- present & Ensemble with gauges            & 0.1$^{\circ}$ \\
\bottomrule
\end{tabular}
\endgroup
\end{table}

\subsection{Known biases against gauges}

Multi-product comparisons of IMERG and its predecessors against
gauge networks have documented three persistent bias modes:
over-detection of light rainfall (the satellite reports rain
when gauges report dry, often associated with sub-pixel cloud
artefacts or shallow drizzle that the satellite retrieves but
the gauge misses), under-detection of heavy events (the
satellite under-estimates extreme rainfall magnitudes, a
limitation of microwave retrieval at the upper tail of
precipitation rates), and a full distributional mismatch
between the satellite and gauge distributions of daily
precipitation~\citep{ref:maggioni2016, ref:tang2020, ref:tan2016}. Over the
tropics specifically, multi-product comparisons report daily
Pearson correlations between $0.3$ and $0.5$ against gauge
networks and reanalyses, with skill improving substantially
only at dekadal or longer aggregation~\citep{ref:ramadhan2022ts, ref:tang2020}. These bias modes motivate the methodological
families reviewed in Section~\ref{sec:lit-methods}; the
empirical magnitudes for IMERG-L over Indonesia are reported in
Chapter~\ref{chap:results}.


% ============================================================
\section{Bias Dimensions and Verification}\label{sec:lit-verification}
% Target: 2-3 pp
% ============================================================

A satellite precipitation product has many ways to be wrong, and
a verification framework that reports a single skill number
hides the dimension in which the product fails. The verification
literature has consequently developed a set of pillar-organised
metrics that target distinct bias dimensions independently, so
that a product's strengths and weaknesses can be diagnosed
rather than averaged together.

\subsection{Four-pillar verification framework}

The four-pillar framework used in this thesis groups verification
metrics into \textit{value adjustment} (mean offset and spread
ratio), \textit{distribution alignment} (cumulative distribution
function comparison and wet-day statistics), \textit{extreme
preservation} (upper-tail percentile ratios and multi-threshold
detection skill), and \textit{event detection} (contingency-table
metrics at a wet-day or heavy-rain threshold). The grouping
follows the conventions in the standard verification references
\citep{ref:wilks2011, ref:ebert2007} and matches the World
Meteorological Organization (WMO) verification standards for
quantitative precipitation forecasts~\citep{ref:wmo1485, ref:wmo1132}. The framework allows
the question \enquote{how well does this product reproduce the
gauge field?} to be decomposed into four narrower questions,
each with a measurable answer.

Pearson correlation between paired daily satellite and gauge
values is reported separately from the four pillars as a
\textit{temporal-skill diagnostic}: it measures the day-by-day
pairing between two series rather than a property of their
marginal distributions. The distinction between
marginal-distribution metrics and temporal-skill metrics becomes
the central object of analysis in
Chapter~\ref{chap:discussion}, and the theoretical basis for
why marginal-correction methods leave Pearson correlation
bounded is developed in Section~\ref{sec:lit-bounds}.

\subsection{Multi-threshold verification}

Categorical skill scores at a single wet-day threshold reduce a
full daily precipitation field to a binary rain-or-no-rain
contingency table and report a single probability of detection,
critical success index, false alarm ratio, and equitable threat
score~\citep{ref:wilks2011}. Single-threshold reporting hides the
threshold-dependence of skill: a product may detect light rain
well and heavy rain poorly, and a single number cannot
distinguish the two regimes. The WMO multi-threshold
verification standard~\citep{ref:wmo1485} reports categorical
skill at a series of thresholds spanning the operational range,
exposing the threshold-dependence directly. The framework adopted
in this study uses seven thresholds spanning this range (1, 5, 10,
20, 50, 100, and 150~mm/day) and reports separate categorical
metrics at each.

\subsection{Station-based out-of-sample validation}

The verification reference may itself contain biases. Gauge
networks have well-documented undercatch behaviour, particularly
for heavy rainfall and at high-elevation or coastal
sites~\citep{ref:adam2003}. Gridded gauge analyses such as the
CPC Unified Gauge-Based Analysis~\citep{ref:chen2008} interpolate
station observations onto a regular grid and inherit the
station-density gradient: cells with many contributing stations
are well-constrained, cells with few stations are smooth-field
reconstructions. A verification that compares the corrected
satellite against the gridded analysis tests the correction's
ability to match a smoothed product, not its ability to match
the original point observations. This distinction motivates the
two-tier validation framework adopted in
Section~\ref{sec:methods-validation}: in-sample comparison
against the gridded reference (the calibration target),
out-of-sample comparison against independent point observations
(the operational user's perspective).


% ============================================================
\section{Bias Correction Method Families}\label{sec:lit-methods}
% Target: 4-5 pp
% ============================================================

The bias correction literature has produced four broad method
families, each targeting one dimension of the distributional
mismatch between a satellite and a reference. The four families
correspond directly to the four stages of the LSEQM+DL pipeline:
linear scaling for the mean, quantile mapping with a parametric
fit for the bulk distribution, extreme-value tail correction for
the upper tail, and deep learning for residual spatial structure.
Each subsection below reviews the principle of one family, its
canonical references, and its acknowledged limits. A unifying
review of the marginal-distribution families is provided by
\citep{ref:maraun2016}.
Table~\ref{tab:biascorr-methods} summarises the principal
strengths and acknowledged weaknesses of the methods considered
in this review, providing a single-page reference against which
the LSEQM+DL pipeline of Chapter~\ref{chap:methods} can be
positioned.

\begin{table}[H]
\centering
\caption[Strengths and weaknesses of bias correction methods considered for daily precipitation]{Strengths and weaknesses of bias correction methods
considered for daily precipitation. The LSEQM+DL pipeline of
Chapter~\ref{chap:methods} combines the first three rows
(LS, EQM, and an extreme-value tail correction) into a sequential
pipeline and adds a deep-learning refinement
stage.\label{tab:biascorr-methods}}
\begingroup\footnotesize
\begin{tabular}{@{}p{3.0cm} p{5.0cm} p{5.0cm}@{}}
\toprule
Method & Strengths & Weaknesses \\
\midrule
Linear Scaling (LS) &
Simple, efficient. Corrects the mean by a single
multiplicative factor. &
Does not address variability or extremes; preserves all
higher-order biases of the satellite distribution. \\
\addlinespace
Empirical Quantile Mapping (EQM) &
Corrects the full marginal distribution (mean, variance,
percentiles) when paired with a parametric fit. &
Computationally demanding; sensitive to the calibration window;
may introduce false alarms at low thresholds when wet-day
matching is misconfigured. \\
\addlinespace
Generalized Pareto Distribution (GPD) tail correction &
Models the upper tail explicitly with extreme-value theory;
corrects heavy-event magnitudes that Gamma EQM under-shoots. &
Requires sufficient exceedances per pixel for a stable fit;
threshold choice is a sensitivity parameter without a single
correct value. \\
\addlinespace
Bias Correction and Spatial Disaggregation (BCSD) &
Combines bias correction with spatial disaggregation; produces
high-resolution corrected fields. &
Computationally expensive; relies on a reference at the
disaggregation target resolution. \\
\addlinespace
Bayesian Model Averaging (BMA) &
Incorporates model-set uncertainty by weighting an ensemble;
improves reliability where multiple candidate models exist. &
Requires multiple models to average; complex implementation;
computationally intensive. \\
\addlinespace
Convolutional Neural Networks (CNN) &
Uses spatial context that per-pixel methods cannot exploit;
refines residual spatial structure of extreme pixels. &
Requires training data; risks amplifying reference uncertainty
in data-sparse regions unless modulated by a confidence weight. \\
\addlinespace
Censored Shifted Mixture Distribution (CSMD) &
Handles zero-inflated daily precipitation directly; performs
well on extremes. &
Complex implementation; computationally intensive; less widely
adopted in the operational literature. \\
\bottomrule
\end{tabular}
\endgroup
\end{table}

\subsection{Linear Scaling}\label{sec:lit-ls}

The simplest family of bias correction methods applies a single
multiplicative or additive factor to the satellite series so
that its mean matches the reference mean over a calibration
window. LS of daily precipitation typically
uses the multiplicative form to preserve non-negativity
\citep{ref:teutschbein2012}:
\begin{equation}
P_{\text{LS}} = P_{\text{sat}} \cdot \frac{\overline{P}_{\text{ref}}}{\overline{P}_{\text{sat}}},
\label{eq:lit-ls}
\end{equation}
where the means are computed over a calibration window matched
to the temporal scale at which the bias is assumed stationary
(monthly, dekadal, or seasonal). LS corrects the first moment
of the marginal distribution by construction but does not change
the distribution shape, the wet-day frequency, or the rank order
of values at a given pixel. As a stand-alone method LS leaves
the higher-order biases of the satellite distribution
unchanged; as the first stage of a multi-stage pipeline it
removes the mean offset so that subsequent stages can work on
the residual distributional mismatch.

\subsection{Quantile Mapping}\label{sec:lit-qm}

Quantile Mapping (QM) corrects the full marginal distribution
by mapping each value in the satellite distribution onto the
corresponding quantile of the reference distribution. The
fundamental relation is the quantile composition
$P_{\text{QM}} = F_{\text{ref}}^{-1}(F_{\text{sat}}(P_{\text{sat}}))$,
where $F_{\text{sat}}$ and $F_{\text{ref}}$ are the cumulative
distribution functions (CDFs) of the satellite and reference
samples respectively. Two main variants of QM differ in how the
CDFs are estimated. Empirical QM uses the empirical CDFs
directly, which is non-parametric and adapts to any
distribution shape but requires a large sample for stable
estimation. Parametric QM fits a chosen distribution family
(typically Gamma for daily precipitation) to each sample,
estimates the parameters by maximum likelihood, and uses the
fitted CDFs for the mapping; parametric QM trades flexibility
for sample-size efficiency.

The Gamma family is the most widely used parametric form for
daily precipitation on wet days, justified by the heavy-tailed
shape and non-negative support of the Gamma distribution and by
empirical fit quality across a wide range of climates
\citep{ref:wilks2011}. The EQM
variant with a Gamma fit, adopted in Stage~2 of the LSEQM+DL
pipeline, traces to the bias correction methodology of
\citep{ref:gudmundsson2012, ref:themessl2012} and is the
standard reference implementation for daily precipitation
correction in climate downscaling contexts.

Quantile mapping methods are known to alter the wet-day
frequency of the satellite series unless the dry-day handling
is corrected explicitly. The wet-day-frequency-matching
correction of \citep{ref:cannon2015} preserves the gauge wet-day
fraction by reclassifying spurious satellite light-rain events
as dry. This correction is the source of the event-detection
trade-off discussed in Chapter~\ref{chap:results}: matching the
dry-day frequency suppresses spurious wet days at the cost of
reduced aggregate probability of detection.

\subsection{Extreme-Value Tail Correction}\label{sec:lit-gpd}

The Gamma distribution fitted in standard EQM is a
two-parameter family whose parameters are dominated by the
bulk of the wet-day sample. The upper tail, the segment
operationally meaningful for flood and drought thresholds, is
fit only as a by-product of forcing the Gamma curve through
the bulk, and is systematically too short for tropical
precipitation. Extreme-value theory provides a more appropriate
parametric form for the tail. The GPD is the canonical extreme-value model in the
peaks-over-threshold framework: under mild regularity
conditions the distribution of exceedances above a high
threshold converges to a GPD with shape and scale parameters
$(\xi, \sigma)$~\citep{ref:coles2001}. The GPD shape parameter
$\xi$ controls tail heaviness, with $\xi > 0$ indicating a
heavy Pareto-like tail typical of tropical convective extremes,
$\xi = 0$ indicating an exponential tail, and $\xi < 0$ indicating
a bounded tail.

The hybrid bias correction strategy is to fit the Gamma family
to the full wet-day sample for the bulk distribution, fit the
GPD to exceedances above a high threshold for the tail, and use
the GPD-based quantile to replace the Gamma-based quantile
above the threshold. The threshold is typically chosen as the
80th to 95th percentile of the wet-day sample, with the trade-off
that higher thresholds isolate the extreme tail more cleanly
but yield fewer exceedances per pixel-dekad and therefore less
stable GPD fits. For the small exceedance samples typical of per-pixel per-dekad
fits, L-moment estimators are often preferred over maximum
likelihood because they are less sensitive to outliers
\citep{ref:hoskingwallis1997}. The present pipeline instead
retains maximum-likelihood estimation with subsample-averaging
to reduce estimator variance (Section~\ref{sec:methods-gpd}).
Stage~3 of the LSEQM+DL pipeline implements this hybrid strategy
with a fixed 80th-percentile threshold.

\subsection{Deep Learning Approaches}\label{sec:lit-dl}

Statistical bias correction methods operate one pixel at a time
and cannot use the spatial context of a precipitation field. A
heavy rainfall event in a single isolated pixel surrounded by
dry pixels is statistically indistinguishable from a coherent
heavy event in a cluster of adjacent pixels under a per-pixel
correction, even though the latter is climatologically more
plausible. Deep learning approaches address this limitation by
operating on the spatial neighbourhood of each pixel.

Convolutional Neural Networks (CNNs) have been applied to
precipitation estimation and statistical downscaling in several
recent studies~\citep{ref:pan2019, ref:banomedina2020}. The principle
is to train a network to map the raw satellite field to the
reference field over a large historical sample, then apply the
trained network at inference time to produce corrected fields.
The network architecture is typically a small convolutional
stack that extracts spatial features at progressively coarser
scales, followed by a reconstruction head that produces a
corrected value at each pixel. Training uses standard
supervised-learning objectives, with mean squared error as the
typical loss for continuous-valued precipitation.

Two design choices have proved important in the application of
CNNs to satellite precipitation correction. First, restricting
the CNN to operate on extreme-pixel values rather than the
full field reduces the function class the network must
approximate and improves sample efficiency in the small-sample
regime typical of per-pixel per-dekad correction. Second,
blending the CNN output with the statistical-correction
intermediate at a fixed weight, rather than replacing it,
preserves the statistical baseline as a fallback in regions
where the CNN has not been adequately trained. Stage~4 of the
LSEQM+DL pipeline implements both of these design choices: the
CNN acts only above the 80th-percentile threshold from Stage~3,
and its output is blended with the LSEQM intermediate at a
nominal weight of $\alpha = 0.70$ for the statistical baseline
and $1 - \alpha = 0.30$ for the CNN.

\subsection{Hybrid Pipelines}\label{sec:lit-hybrid}

The four method families described above target complementary
bias dimensions, and combining them as sequential stages is a
common pattern in the bias correction literature
\citep{ref:maraun2016}. A typical hybrid pipeline applies LS for
the mean, EQM for the bulk distribution, and either GPD or a
separate heavy-rain correction for the tail; recent work has
added deep-learning refinement as a final stage. The pipelines
differ in the parametric family used for EQM (Gamma, Weibull,
or non-parametric), the threshold choice for GPD, and the
deep-learning architecture, but the underlying logic of
matching one stage to one bias dimension is consistent across
the literature. The LSEQM+DL pipeline described in
Chapter~\ref{chap:methods} follows this pattern with a specific
implementation of each stage and an explicit
station-density-aware blending design that modulates the deep
learning contribution by local gauge density.


% ============================================================
\section{Theoretical Bounds on Marginal Correction}\label{sec:lit-bounds}
% ============================================================

The bias correction families reviewed in
Section~\ref{sec:lit-methods} share a structural property that
the verification literature has examined closely: they correct
the marginal distribution of a series at each location, and they
do so through transformations that preserve the rank order of
the values they act on. This property has consequences for what
such methods can and cannot improve, and those consequences
frame the central empirical result of the present study. This
section reviews the two established results that together bound
the time-dependent skill attainable by any marginal correction.

The first is rank preservation. Quantile mapping, in all of its
parametric and non-parametric variants, maps each value to the
reference quantile that matches its rank in the source
distribution. When the mapping is strictly monotonic, which
holds for empirical quantile mapping with a continuous reference
CDF and for parametric mapping through a Gamma or Generalized
Pareto fit, the rank order of the values at a location is
preserved exactly: if the source series reports a higher value
on one day than on another, the corrected series does too
\citep{ref:cannon2015}. Cannon and colleagues examined this
property in the context of trend preservation and showed that
standard quantile mapping can distort climate-change trends
precisely because it operates on ranks rather than on the
temporal sequence; the same rank-preservation mechanism that
distorts trends also fixes the day-by-day ordering of a
corrected series relative to its source.

The second is variance inflation without time-dependent skill
gain. \citet{ref:maraun2013} analysed quantile mapping as
applied to regional climate model output and showed that the
method inflates or deflates the variance of a series to match
the reference marginal, but that this spread adjustment does not
add information about the timing of individual events. The
adjustment changes the amplitude of the series without changing
which days are wet and which are dry relative to the source
ordering. Spread alignment and day-by-day pairing are, in this
analysis, independent properties of a paired series: a
correction can move one without moving the other.

Taken together, the two results imply a bound that is directly
relevant to satellite precipitation correction. The Pearson
correlation between a corrected satellite series and an external
gauge series measures the day-by-day pairing of the two series.
If a marginal correction preserves the rank order of the
satellite series, it cannot change which satellite day is paired
with which gauge day, and the correlation is therefore bounded
above by the correlation between the raw satellite series and the
gauge. The amplitude correction that moves the standard deviation
ratio toward unity has no corresponding effect on the
correlation, because correlation is invariant under the monotone
rescaling that quantile mapping applies. This is not a deficiency
of any particular implementation; it is a property of the
marginal correction family as a whole, and any method built from
linear scaling, quantile mapping, and tail adjustment inherits
it.

The empirical magnitude of the bound for daily IMERG over the
tropics is well documented. Multi-product comparisons of IMERG
against gauge networks and reanalyses
report daily correlations in the range of $0.3$ to $0.5$ across
products and seasons \citep{ref:tang2020}, and evaluations over
the Indonesian Maritime Continent report daily correlations of
similar magnitude that rise substantially only when the
aggregation window is widened to dekadal or monthly scales
\citep{ref:ramadhan2022ts}. These values define the empirical
ceiling within which any corrected daily product over the region
must operate. The present study tests the rank-preservation
bound directly by reporting the Pearson correlation at each
correction stage against independent gauges
(Chapter~\ref{chap:results}) and develops its structural
explanation as the central argument of
Chapter~\ref{chap:discussion}.


% ============================================================
\section{Reproducible Computational Notebook Workflows in Climate and Earth Observation Analytics}\label{sec:lit-reproducibility}
% ============================================================

The empirical claim of a bias correction study rests on the
quality of the calibration data and the verification framework,
but it also rests on the ability of an independent reader to
re-execute the workflow and arrive at the same numbers. The
reproducibility discourse in computational science
\citep{ref:peng2011, ref:stodden2014} and the FAIR (Findable,
Accessible, Interoperable, Reusable) principles for scientific
data management \citep{ref:wilkinson2016} have
made independent verifiability a methodological expectation in
the broader literature. The bias correction literature has
adopted these expectations unevenly: many influential studies
publish methods without releasing the corresponding code or
data, leaving any claim of reproduction to a reader who must
re-implement the method from the paper alone. The present study
adopts an explicit reproducibility design that is reviewed in
Section~\ref{sec:methods-reproducibility} and demonstrated
empirically in Section~\ref{sec:results-reproducibility}; the
present section reviews the supporting literature on
computational notebook workflows that the design draws on.

\subsection{Computational notebooks in climate and Earth observation analysis}

Computational notebooks have become a standard medium for
documenting and sharing analytical workflows in climate science
and Earth observation \cite{ref:perkel2018, ref:wagemann2022}. A notebook integrates code, results, and
narrative in a single document, making the analytical process
transparent and re-runnable on the same data. Jupyter is the
most widely used implementation in the scientific Python
ecosystem and is the format adopted in the present study; the
Jupyter notebook format is portable across execution
environments, opening in either a local installation or a
cloud-hosted runtime such as Google Colab with no change to
the underlying file. For climate science specifically,
computational notebooks have been used to integrate satellite
imagery, station records, and model outputs into single
self-documenting workflows \citep{ref:gentemann2021, ref:kluyver2016},
and to teach climate-data analytics in a way that makes the
underlying processing steps explicit and modifiable
\citep{ref:wagemann2022}. The notebook format
preserves the connection between the prose claim and the code
that produced it, which is precisely the connection that a
peer-reviewed paper without released code cannot guarantee.

The integration of computational notebooks with cloud-based
geospatial platforms has further lowered the barrier to running
a published workflow. Google Earth Engine and Microsoft
Planetary Computer expose their data catalogues through
notebook-friendly Python libraries such as \texttt{geemap},
allowing a reader to execute a published analysis in a browser
without local installation \cite{ref:gorelick2017, ref:wu2020}. Google Colab itself
provides a no-installation execution environment for the same
notebook files that run locally under Jupyter, with the cloud
runtime additionally providing free Graphics Processing Unit (GPU)
access for the deep learning step of the present study. These platforms make
geospatial and machine-learning computation accessible to
readers without high-performance local infrastructure and have
been argued to democratise access to the satellite data
archives the bias correction literature draws on
\citep{ref:gentemann2021, ref:yin2019}. The combination of
computational notebooks and cloud platforms is also a teaching
aid: \citep{ref:haedrich2023} use notebook-based tutorials in the
Geographic Resources Analysis Support System (GRASS) Geographic
Information System (GIS) to teach advanced geospatial modelling,
and \citep{ref:boeing2021} argue that the notebook format makes
geographic analysis more nimble and extensible than conventional
Graphical User Interface (GUI)-driven workflows.

For the reproducibility design adopted in this thesis, two
properties of the computational notebook format are decisive.
First, the six numbered notebooks of the LSEQM+DL workflow
expose every correction stage and every verification step as a
runnable artefact tied to specific functions in the source
codebase, so a reader following the documentation can move
directly to the runnable cell and back \cite{ref:perkel2018,
ref:granger2021, ref:beg2021}. Second, the notebooks can be opened in either
a local Jupyter installation or a cloud-hosted Google Colab
session, which means the workflow is portable across the local
development environment the author used and the browser-hosted
environment a reader is likely to start with. The empirical
demonstration of this portability is reported in
Section~\ref{sec:results-reproducibility}.

\subsection{Notebook-based workflows in bias correction}

Within the bias correction literature specifically,
notebook-based release remains the exception rather than the
rule. A reader who wishes to compare the LSEQM+DL pipeline of
Chapter~\ref{chap:methods} against a published alternative is
typically able to find the method description but not the code.
This gap between method documentation and runnable artefact is
the methodological position that the present study takes a
deliberate stance against: the chapter argues that the empirical
claim of a bias correction study should be re-executable by an
independent reader, and the implementation reported in
Section~\ref{sec:methods-reproducibility} is structured to make
that claim verifiable rather than aspirational. The wider
literature on open scientific workflows
\citep{ref:wilkinson2016, ref:giuliani2019} supports this stance,
though it has not yet become a community norm in the bias
correction sub-literature.


% ============================================================
\section{Maritime Continent Precipitation Context}\label{sec:lit-maritime-continent}
% Target: 2-3 pp
% Status: SKELETON - drafted in LitRev v2
% ============================================================

The Indonesian Maritime Continent presents a precipitation
regime that is among the most challenging in the world for
satellite estimation, and several of its features bear directly
on the verification and correction problem this study addresses.

Rainfall over the archipelago is dominated by deep convection
organised by the interaction between the large-scale monsoon
circulation and intense local land-sea breeze systems
\citep{ref:aldrian2003, ref:wu2008}. Two monsoon phases set the seasonal
cycle: the northwest monsoon (roughly November to March) brings
the primary wet season to most of the archipelago, while the
southeast monsoon (roughly May to September) brings the dry
season to the south and a secondary regime to the north. The
seasonal migration of the Inter-Tropical Convergence Zone and
the influence of the Madden-Julian Oscillation add intraseasonal
variability on top of the monsoon cycle. The result is a daily
precipitation field with a heavy convective tail and strong
spatial heterogeneity, which is the regime the Generalized
Pareto tail adjustment and the station-density-aware blending of
Chapter~\ref{chap:methods} are designed to handle.

The diurnal cycle is the feature most consequential for daily
verification. Convection over the large islands peaks in the
local afternoon, typically between 14:00 and 17:00 local time,
driven by daytime heating and the sea-breeze convergence over
land \citep{ref:wu2008}. Indonesia spans three standard
time zones, from UTC+7 in the west (Sumatra, Java) through UTC+8
(Bali, Kalimantan, Sulawesi) to UTC+9 in the east (Maluku,
Papua). The IMERG product is accumulated on a UTC calendar day,
so the local convective afternoon falls in the UTC 07:00 to
10:00 band for the western archipelago and earlier in UTC terms
for the east. A daily aggregation window aligned to UTC
therefore splits the local convective afternoon across two
calendar days for stations away from the UTC meridian, which
degrades the day-by-day pairing between a UTC-day satellite
total and a local-day gauge total by construction. This calendar
artefact is one contributor to the daily correlation ceiling
discussed in Section~\ref{sec:lit-bounds}, and it motivates the
local-day re-aggregation path that Chapter~\ref{chap:discussion}
identifies as the most actionable next step for the region.

The gauge network that anchors both the CPC-UNI reference and
the independent validation is strongly heterogeneous in space.
The Indonesian Meteorological, Climatological, and Geophysical Agency (BMKG) station network is dense over Java and southern
Sumatra, where population and agricultural value are
concentrated, and progressively sparser eastward into Maluku,
Nusa Tenggara, and Papua \citep{ref:chen2008}. Because CPC-UNI is
a gauge-interpolated analysis, its quality follows the same
gradient: the reference is well constrained in the west and
weakly constrained in the east. This spatial gradient in
reference quality is the reason the present study introduces a
station-density-aware blending design rather than applying a
uniform deep-learning correction across the domain
(Section~\ref{sec:methods-blending}); it lets the framework
revert to the purely statistical correction where the reference
cannot support a learned refinement.

Geostationary observation provides the diurnal-cycle reference
that any future local-day re-aggregation would need. The
Himawari-9 platform covers the Maritime Continent at high
temporal cadence (on the order of ten minutes) and
kilometre-scale spatial resolution, which is sufficient to track
the convective afternoon and to align a re-aggregated daily
window to the local cycle. The present study does not perform
this re-aggregation, but the regional context established here
frames it as the natural extension that
Chapter~\ref{chap:discussion} proposes.


% ============================================================
\section*{Chapter summary}
% Status: SKELETON - drafted in LitRev v2 once all six sections are filled
% ============================================================

This review has established the foundations on which the present
study builds. The satellite precipitation products of
Section~\ref{sec:lit-imerg} define the input to be corrected and
its known error structure; the bias dimensions and verification
standards of Section~\ref{sec:lit-verification} define the
four-pillar framework against which the correction is judged; and
the bias correction families of Section~\ref{sec:lit-methods}
define the four stages that the LSEQM+DL pipeline composes. The
theoretical bounds of Section~\ref{sec:lit-bounds} establish, in
advance of any result, what a marginal correction can and cannot
move: it can align the marginal distribution but cannot lift the
day-by-day correlation above the ceiling set by the raw
satellite pairing. The reproducibility literature of
Section~\ref{sec:lit-reproducibility} establishes the
methodological expectation of independent verifiability that the
study's open pipeline is designed to meet, and the regional
context of Section~\ref{sec:lit-maritime-continent} establishes
the convective regime, the UTC calendar artefact, and the gauge
density gradient that shape both the correction design and its
interpretation. Chapter~\ref{chap:methods} now describes the
pipeline, the blending design, the reproducible implementation,
and the verification framework in full.
